\documentclass[12pt, a4paper]{article}

% ── Codificación y tipografía ──────────────────────────────────────────────────
\usepackage[utf8]{inputenc}
\usepackage[T1]{fontenc}
\usepackage[spanish]{babel}
\usepackage{lmodern}
\usepackage{microtype}

% ── Geometría ─────────────────────────────────────────────────────────────────
\usepackage[
  top=2.8cm, bottom=2.8cm,
  left=3.2cm, right=3.2cm,
  headheight=15pt
]{geometry}

% ── Color y gráficos ──────────────────────────────────────────────────────────
\usepackage[dvipsnames, table]{xcolor}
\definecolor{primary}{HTML}{1A237E}
\definecolor{secondary}{HTML}{1565C0}
\definecolor{accent}{HTML}{0277BD}
\definecolor{codebg}{HTML}{F5F5F5}
\definecolor{tabheader}{HTML}{1565C0}
\definecolor{rowalt}{HTML}{EEF4FB}
\definecolor{rulecolor}{HTML}{90CAF9}

% ── Encabezado y pie de página ────────────────────────────────────────────────
\usepackage{fancyhdr}
\pagestyle{fancy}
\fancyhf{}
\fancyhead[L]{\small\textcolor{secondary}{Equidad Algorítmica – SaberXAI}}
\fancyhead[R]{\small\textcolor{secondary}{Documentación del Dataset}}
\fancyfoot[C]{\small\textcolor{gray}{\thepage}}
\renewcommand{\headrulewidth}{0.4pt}
\renewcommand{\footrulewidth}{0pt}

% ── Secciones con color ───────────────────────────────────────────────────────
\usepackage{titlesec}
\titleformat{\section}
  {\Large\bfseries\color{primary}}
  {\thesection.}{0.6em}{}[{\color{rulecolor}\titlerule[0.8pt]}]
\titleformat{\subsection}
  {\large\bfseries\color{secondary}}
  {\thesubsection.}{0.5em}{}
\titleformat{\subsubsection}
  {\normalsize\bfseries\color{accent}}
  {\thesubsubsection.}{0.4em}{}
\titlespacing{\section}{0pt}{18pt}{8pt}
\titlespacing{\subsection}{0pt}{12pt}{4pt}
\titlespacing{\subsubsection}{0pt}{8pt}{3pt}

% ── Tablas ────────────────────────────────────────────────────────────────────
\usepackage{booktabs}
\usepackage{tabularx}
\usepackage{array}
\usepackage{multirow}
\renewcommand{\arraystretch}{1.35}

% ── Código ────────────────────────────────────────────────────────────────────
\usepackage{listings}
\lstset{
  backgroundcolor=\color{codebg},
  basicstyle=\ttfamily\small,
  breaklines=true,
  frame=single,
  framerule=0.4pt,
  rulecolor=\color{gray!40},
  xleftmargin=8pt,
  xrightmargin=8pt,
  aboveskip=8pt,
  belowskip=6pt,
  columns=fullflexible,
  keepspaces=true,
  showstringspaces=false,
  language=Python,
  keywordstyle=\color{secondary}\bfseries,
  stringstyle=\color{accent},
  commentstyle=\color{gray}\itshape,
}

% ── Hiperenlaces ──────────────────────────────────────────────────────────────
\usepackage[
  colorlinks=true,
  linkcolor=secondary,
  urlcolor=accent,
  citecolor=secondary,
  pdftitle={Documentación del Dataset ICFES Saber 11},
  pdfauthor={Jacobo Ayala},
  pdfsubject={Decisiones de preparación del dataset para análisis de equidad}
]{hyperref}

% ── Utilidades ────────────────────────────────────────────────────────────────
\usepackage{enumitem}
\setlist[itemize]{leftmargin=1.4em, itemsep=2pt, topsep=4pt}
\setlist[enumerate]{leftmargin=1.4em, itemsep=2pt, topsep=4pt}
\usepackage{parskip}
\setlength{\parskip}{5pt}
\usepackage{caption}
\captionsetup{font=small, labelfont=bf, labelsep=period}
\usepackage{tcolorbox}
\tcbuselibrary{skins, breakable}

\newtcolorbox{nota}{
  enhanced, breakable,
  colback=rowalt,
  colframe=secondary,
  boxrule=0.6pt,
  left=8pt, right=8pt, top=5pt, bottom=5pt,
  fonttitle=\bfseries\small,
  title=Nota,
}
\newtcolorbox{decision}{
  enhanced, breakable,
  colback=codebg,
  colframe=accent,
  boxrule=0.5pt,
  left=8pt, right=8pt, top=5pt, bottom=5pt,
}

% ── Comando auxiliar para celdas de encabezado de tabla ───────────────────────
\newcommand{\thead}[1]{\cellcolor{tabheader}\textcolor{white}{\textbf{#1}}}

% ══════════════════════════════════════════════════════════════════════════════
\begin{document}

% ── Portada ───────────────────────────────────────────────────────────────────
\begin{titlepage}
  \centering
  \vspace*{3cm}
  {\color{primary}\rule{\linewidth}{1.5pt}}\vspace{0.6cm}

  {\Huge\bfseries\color{primary} Documentación del Dataset\\[0.3em]
  ICFES Saber 11}\vspace{0.5cm}

  {\color{primary}\rule{\linewidth}{1.5pt}}\vspace{1.2cm}

  {\large\color{secondary} Proyecto: Equidad Algorítmica — SaberXAI}\vspace{0.4cm}

  {\normalsize Predicción de puntaje global con deep learning\\
  e interpretabilidad orientada a brechas de equidad}\vspace{2cm}

  \begin{tabular}{rl}
    \textbf{Autor:}    & Jacobo Ayala \\[4pt]
    \textbf{Periodos:} & 2014-2 a 2024-2 (21 semestres) \\[4pt]
    \textbf{Fuente:}   & ICFES — Examen Saber 11 \\[4pt]
    \textbf{Dataset:}  & \texttt{jacoboayala/datos-icfes-11} (Kaggle) \\[4pt]
    \textbf{Fecha:}    & \today \\
  \end{tabular}

  \vfill
  {\small\color{gray}
  Este documento consolida las decisiones técnicas y analíticas tomadas
  durante la construcción y limpieza del dataset. Cubre la selección de
  variables, la exclusión de periodos, las transformaciones aplicadas y las
  justificaciones metodológicas de cada elección.}
\end{titlepage}

% ── Tabla de contenidos ───────────────────────────────────────────────────────
\tableofcontents
\newpage

% ══════════════════════════════════════════════════════════════════════════════
\section{Contexto del proyecto}
% ══════════════════════════════════════════════════════════════════════════════

Este proyecto aplica deep learning sobre el dataset ICFES Saber 11 para predecir
\texttt{punt\_global} como variable objetivo. El objetivo central es dar
\textbf{interpretabilidad al modelo} para analizar brechas de equidad a lo largo
de cinco dimensiones:

\begin{itemize}
  \item Área geográfica (urbano / rural)
  \item Estrato socioeconómico
  \item Género del estudiante
  \item Tipo y naturaleza del colegio
  \item Departamento de ubicación
\end{itemize}

La selección de variables y todas las decisiones de preparación del dataset están
guiadas por tres principios:

\begin{enumerate}
  \item Eliminar \textit{data leakage} que haría el modelo trivial.
  \item Conservar únicamente \textit{features} que expliquen condiciones del
        estudiante, no sus resultados.
  \item Garantizar que los métodos de interpretabilidad (SHAP u otros) reflejen
        factores reales de equidad.
\end{enumerate}

% ══════════════════════════════════════════════════════════════════════════════
\section{Estructura del dataset longitudinal}
% ══════════════════════════════════════════════════════════════════════════════

El dataset consolidado cubre \textbf{21 periodos semestrales} (2014-2 a 2024-2).
Cada periodo proviene de un archivo \texttt{.txt} separado con delimitador
\texttt{;} publicado por el ICFES.

\subsection{Variación de esquemas entre periodos}

Los archivos no comparten un esquema homogéneo. Se identificaron \textbf{10
esquemas distintos}; la siguiente tabla resume los grupos principales:

\vspace{4pt}
\begin{tabularx}{\textwidth}{>{\ttfamily}l c X}
  \toprule
  \thead{Grupo de archivos} & \thead{Columnas} & \thead{Periodos} \\
  \midrule
  \rowcolor{rowalt}
  20142                              & 57 & Semestre de referencia mínimo \\
  20151                              & 58 & +\texttt{fami\_numlibros} \\
  \rowcolor{rowalt}
  20152                              & 59 & +\texttt{estu\_tieneetnia} \\
  20161–20162                        & 72 & +13 columnas de desempeño y percentiles \\
  \rowcolor{rowalt}
  20171–20181, 20182–20201, 20212    & 85 & +13 columnas de hábitos y familia \\
  20202, 20211, 20221–20242          & 83–84 & Esquema estabilizado reciente \\
  \rowcolor{rowalt}
  \textbf{20141}                     & \textbf{138} & \textbf{Excluido (ver \S\ref{sec:exclusion})} \\
  \bottomrule
\end{tabularx}

% ══════════════════════════════════════════════════════════════════════════════
\section{Exclusión del periodo 2014-1}
\label{sec:exclusion}
% ══════════════════════════════════════════════════════════════════════════════

\subsection{Decisión}

El archivo \texttt{Examen\_Saber\_11\_20141.txt} se excluye del análisis principal.

\subsection{Razón técnica}

El periodo 2014-1 tiene una estructura de columnas significativamente distinta
al resto:

\begin{itemize}
  \item \textbf{2014-1}: 138 columnas (incluye variables vocacionales, de
        trayectoria escolar y puntajes por módulo del esquema anterior al
        rediseño ICFES de 2014).
  \item \textbf{Demás periodos}: entre 57 y 85 columnas.
\end{itemize}

El impacto en la intersección de columnas comunes es determinante:

\vspace{4pt}
\begin{tabular}{lcc}
  \toprule
  \thead{Escenario} & \thead{Columnas en la intersección} \\
  \midrule
  \rowcolor{rowalt}
  Incluyendo 2014-1 & 46 \\
  Excluyendo 2014-1 & \textbf{55} \\
  \bottomrule
\end{tabular}
\vspace{6pt}

Excluir 2014-1 mejora la consistencia del esquema longitudinal y reduce la pérdida
de información útil en variables comparables entre periodos.

\subsection{Implementación}

La exclusión está codificada en \texttt{merge\_columnas\_comunes.py}:

\begin{lstlisting}
def listar_archivos_datos(carpeta_datos: Path) -> Iterable[Path]:
    archivos = sorted(carpeta_datos.glob("Examen_Saber_11_*.txt"))
    return [a for a in archivos if "20141" not in a.stem]
\end{lstlisting}

% ══════════════════════════════════════════════════════════════════════════════
\section{Construcción del dataset base (55 columnas)}
% ══════════════════════════════════════════════════════════════════════════════

\subsection{Lectura con \texttt{dtype=str}}

Todos los archivos crudos se leen forzando \texttt{dtype=str}:

\begin{lstlisting}
df = pd.read_csv(archivo, sep=";", dtype=str, encoding="utf-8")
\end{lstlisting}

Los 21 periodos presentan esquemas distintos; \texttt{pandas} inferiría tipos de
forma inconsistente. Por ejemplo, un campo que en 2014-2 contiene sólo valores
numéricos se leería como \texttt{int64}, pero en 2020-1 podría incluir \texttt{"N/A"}
y cambiaría a \texttt{object}. Forzar \texttt{dtype=str} garantiza que la
concatenación no falle ni produzca coerciones silenciosas. La conversión a tipos
definitivos ocurre en el notebook de presentación, después del merge.

\subsection{Columnas faltantes en periodos antiguos}

Cuando un periodo no contiene una columna de las 55 comunes, se rellena con
\texttt{pd.NA} en lugar de omitir el archivo:

\begin{lstlisting}
faltantes = [c for c in columnas_objetivo if c not in df.columns]
for col in faltantes:
    df[col] = pd.NA
\end{lstlisting}

Esto preserva la estructura uniforme y permite el \texttt{concat} sin pérdida de
filas.

\subsection{Limpieza de nulos}

Sobre el consolidado se aplica eliminación de filas con al menos un valor nulo:

\vspace{4pt}
\begin{tabular}{lr}
  \toprule
  \thead{Métrica} & \thead{Valor} \\
  \midrule
  \rowcolor{rowalt}
  Filas originales                 & 7\,239\,505 \\
  Filas conservadas (sin nulos)    & 4\,821\,037 \\
  \rowcolor{rowalt}
  Filas excluidas                  & 2\,418\,468 \\
  Porcentaje conservado            & 66.59\% \\
  \bottomrule
\end{tabular}

% ══════════════════════════════════════════════════════════════════════════════
\section{Selección de variables}
\label{sec:features}
% ══════════════════════════════════════════════════════════════════════════════

Se parte de las 55 columnas comunes y se eliminan 29, quedando 28 variables más
el target (\texttt{punt\_global}).

\subsection{Variables eliminadas}

\subsubsection{Data leakage — sub-puntajes}

\vspace{4pt}
\begin{tabularx}{\textwidth}{>{\ttfamily}l X}
  \toprule
  \thead{Variable} & \thead{Razón} \\
  \midrule
  \rowcolor{rowalt}
  punt\_matematicas          & Componente directo de \texttt{punt\_global} \\
  punt\_lectura\_critica     & Componente directo de \texttt{punt\_global} \\
  \rowcolor{rowalt}
  punt\_c\_naturales         & Componente directo de \texttt{punt\_global} \\
  punt\_sociales\_ciudadanas & Componente directo de \texttt{punt\_global} \\
  \rowcolor{rowalt}
  punt\_ingles               & Componente directo de \texttt{punt\_global} \\
  desemp\_ingles             & Nivel de desempeño derivado de \texttt{punt\_ingles} \\
  \bottomrule
\end{tabularx}
\vspace{6pt}

\texttt{punt\_global} se calcula como función directa de estos sub-puntajes.
Incluirlos como \textit{features} haría que el modelo simplemente reconstruyera
la fórmula del ICFES, y los métodos de interpretabilidad mostrarían que
``matemáticas importa'' por construcción matemática, no por hallazgos de equidad.

\subsubsection{Identificadores administrativos}

\vspace{4pt}
\begin{tabularx}{\textwidth}{>{\ttfamily}l X}
  \toprule
  \thead{Variable} & \thead{Razón} \\
  \midrule
  \rowcolor{rowalt}
  estu\_estudiante                & Número de documento — identificador puro \\
  estu\_consecutivo               & Número secuencial ICFES por periodo \\
  \rowcolor{rowalt}
  cole\_cod\_dane\_establecimiento & Código DANE — sin valor predictivo \\
  cole\_cod\_dane\_sede           & Código DANE sede \\
  \rowcolor{rowalt}
  cole\_cod\_depto\_ubicacion     & Redundante con \texttt{cole\_depto\_ubicacion} \\
  cole\_cod\_mcpio\_ubicacion     & Redundante con \texttt{cole\_mcpio\_ubicacion} \\
  \rowcolor{rowalt}
  cole\_codigo\_icfes             & Código interno ICFES \\
  estu\_cod\_depto\_presentacion  & Redundante con \texttt{estu\_depto\_presentacion} \\
  \rowcolor{rowalt}
  estu\_cod\_mcpio\_presentacion  & Redundante con \texttt{estu\_mcpio\_presentacion} \\
  estu\_cod\_reside\_depto        & Redundante con \texttt{estu\_depto\_reside} \\
  \rowcolor{rowalt}
  estu\_cod\_reside\_mcpio        & Redundante con \texttt{estu\_mcpio\_reside} \\
  \bottomrule
\end{tabularx}

\subsubsection{Nombres de alta cardinalidad}

\vspace{4pt}
\begin{tabularx}{\textwidth}{>{\ttfamily}l X}
  \toprule
  \thead{Variable} & \thead{Razón} \\
  \midrule
  \rowcolor{rowalt}
  cole\_nombre\_establecimiento & Texto libre con miles de valores únicos \\
  cole\_nombre\_sede            & Ídem; la información geográfica relevante
                                  ya está en \texttt{cole\_depto\_ubicacion} \\
  \bottomrule
\end{tabularx}

\subsubsection{Redundantes con columnas conservadas}

\vspace{4pt}
\begin{tabularx}{\textwidth}{>{\ttfamily}l l X}
  \toprule
  \thead{Variable} & \thead{Redundante con} & \thead{Razón} \\
  \midrule
  \rowcolor{rowalt}
  periodo                  & \texttt{anio + semestre} & Se descompuso en variables más útiles \\
  estu\_depto\_presentacion & \texttt{cole\_depto\_ubicacion} & El lugar de examen coincide con el colegio en la mayoría de casos \\
  \rowcolor{rowalt}
  estu\_mcpio\_presentacion & \texttt{cole\_depto\_ubicacion} & Mayor cardinalidad sin ganancia informativa \\
  estu\_mcpio\_reside       & \texttt{estu\_depto\_reside} & El departamento es suficiente para capturar geografía de residencia \\
  \bottomrule
\end{tabularx}

\subsubsection{Varianza casi nula}

\vspace{4pt}
\begin{tabularx}{\textwidth}{>{\ttfamily}l l X}
  \toprule
  \thead{Variable} & \thead{Distribución} & \thead{Razón} \\
  \midrule
  \rowcolor{rowalt}
  estu\_privado\_libertad & 100\% N & Sin varianza útil \\
  estu\_pais\_reside      & 99.57\% COLOMBIA & 97 valores únicos pero concentrado en uno \\
  \rowcolor{rowalt}
  estu\_nacionalidad      & 99.43\% COLOMBIA & Ídem \\
  \bottomrule
\end{tabularx}

\subsubsection{Índices socioeconómicos compuestos}

\vspace{4pt}
\begin{tabularx}{\textwidth}{>{\ttfamily}l X}
  \toprule
  \thead{Variable} & \thead{Razón} \\
  \midrule
  \rowcolor{rowalt}
  estu\_inse\_individual & INSE calculado por ICFES a partir de variables \texttt{fami\_*}. Redundante si se conservan las variables base. \\
  estu\_nse\_individual  & NSE individual — mismo argumento. \\
  \bottomrule
\end{tabularx}
\vspace{6pt}

Se prefieren las variables \texttt{fami\_*} desagregadas porque son más
interpretables para el análisis de equidad: permiten identificar qué dimensión
específica del nivel socioeconómico (educación de padres, bienes del hogar,
estrato) tiene mayor impacto.

\subsubsection{Flag administrativo derivado}

\vspace{4pt}
\begin{tabularx}{\textwidth}{>{\ttfamily}l X}
  \toprule
  \thead{Variable} & \thead{Razón} \\
  \midrule
  \rowcolor{rowalt}
  estu\_agregado & Indica si el estudiante entra en estadísticas agregadas del colegio. Flag derivado de condiciones ya capturadas por otras variables. No es una causa del desempeño. \\
  \bottomrule
\end{tabularx}

\subsection{Resumen de eliminaciones}

\vspace{4pt}
\begin{tabular}{lc}
  \toprule
  \thead{Categoría} & \thead{Columnas eliminadas} \\
  \midrule
  \rowcolor{rowalt}
  Data leakage                    & 6  \\
  Identificadores administrativos & 11 \\
  \rowcolor{rowalt}
  Alta cardinalidad               & 2  \\
  Redundantes                     & 4  \\
  \rowcolor{rowalt}
  Baja varianza                   & 3  \\
  Compuestos redundantes          & 2  \\
  \rowcolor{rowalt}
  Flags administrativos           & 1  \\
  \midrule
  \textbf{Total eliminadas}       & \textbf{29} \\
  \textbf{Total conservadas}      & \textbf{28} \\
  \bottomrule
\end{tabular}

\subsection{Variables conservadas (28)}

\vspace{4pt}
\begin{tabularx}{\textwidth}{l X}
  \toprule
  \thead{Grupo} & \thead{Variables} \\
  \midrule
  \rowcolor{rowalt}
  Colegio     & \texttt{cole\_area\_ubicacion}, \texttt{cole\_bilingue},
                \texttt{cole\_calendario}, \texttt{cole\_caracter},
                \texttt{cole\_depto\_ubicacion}, \texttt{cole\_genero},
                \texttt{cole\_jornada}, \texttt{cole\_mcpio\_ubicacion},
                \texttt{cole\_naturaleza}, \texttt{cole\_sede\_principal} \\
  Estudiante  & \texttt{estu\_depto\_reside}, \texttt{estu\_fechanacimiento},
                \texttt{estu\_genero}, \texttt{estu\_grado},
                \texttt{estu\_tipodocumento} \\
  \rowcolor{rowalt}
  Familia     & \texttt{fami\_estratovivienda}, \texttt{fami\_educacionmadre},
                \texttt{fami\_educacionpadre}, \texttt{fami\_cuartoshogar},
                \texttt{fami\_personashogar}, \texttt{fami\_tienecomputador},
                \texttt{fami\_tieneinternet}, \texttt{fami\_tienelavadora},
                \texttt{fami\_tieneautomovil}, \texttt{fami\_tieneserviciotv} \\
  Temporal    & \texttt{anio}, \texttt{semestre} \\
  \rowcolor{rowalt}
  Target      & \texttt{punt\_global} \\
  \bottomrule
\end{tabularx}

\subsection{Notas sobre variables con justificación no obvia}

\subsubsection{\texttt{estu\_grado}}

Se conserva a pesar de que el 88\% de los registros corresponden a grado 11.
El 11.91\% restante corresponde al \textbf{grado 26 (validantes de bachillerato)}:
adultos que certifican su bachillerato a través del examen sin haber cursado
grado 11 en una institución. Esta población tiene características socioeconómicas
y de desempeño sistemáticamente distintas, relevante para el análisis de equidad.

\subsubsection{\texttt{cole\_mcpio\_ubicacion}}

Aporta granularidad geográfica dentro de cada departamento, relevante para
análisis de brechas entre municipios capitales y no capitales.
\texttt{cole\_depto\_ubicacion} sola no captura esta variación intra-departamental.
Cardinalidad tras normalización: 1\,081 valores únicos; se asume que los embeddings
del modelo de deep learning la manejarán adecuadamente.

\subsubsection{\texttt{estu\_tipodocumento}}

Actúa como proxy del estatus migratorio del estudiante (CC, TI, CE, pasaporte,
etc.). Estudiantes con Cédula de Extranjería o pasaporte representan una población
con condiciones de acceso diferenciadas, relevante para el análisis de equidad.
Cardinalidad: menos de 10 categorías; no genera presión en los embeddings.

% ══════════════════════════════════════════════════════════════════════════════
\section{Transformaciones aplicadas al dataset}
% ══════════════════════════════════════════════════════════════════════════════

\subsection{Normalización geográfica con \texttt{strip\_accents}}

El ICFES registra nombres con y sin tilde según el periodo
(e.g., \emph{Bogotá} vs.\ \emph{Bogota}, \emph{Córdoba} vs.\ \emph{Cordoba}).
Se aplica normalización Unicode NFKD seguida de conversión a ASCII y
\texttt{.title()}:

\begin{lstlisting}
def strip_accents(s):
    return unicodedata.normalize("NFKD", str(s)) \
        .encode("ascii", "ignore").decode("ascii").strip().title()

GEO_COLS = ["cole_depto_ubicacion", "cole_mcpio_ubicacion", "estu_depto_reside"]
for col in GEO_COLS:
    df[col] = df[col].map(strip_accents)
\end{lstlisting}

\textbf{Impacto medido:}

\vspace{4pt}
\begin{tabular}{lrrr}
  \toprule
  \thead{Columna} & \thead{Antes} & \thead{Después} & \thead{Reducción} \\
  \midrule
  \rowcolor{rowalt}
  \texttt{cole\_depto\_ubicacion} & 34 & 33 & $-1$ \\
  \texttt{cole\_mcpio\_ubicacion} & 1\,396 & 1\,081 & $-315$ \\
  \rowcolor{rowalt}
  \texttt{estu\_depto\_reside}    & 35 & 34 & $-1$ \\
  \bottomrule
\end{tabular}
\vspace{6pt}

Sin esta normalización el modelo aprendería representaciones distintas para el
mismo lugar, contaminando los \textit{embeddings} geográficos y los análisis de
brecha departamental.

\subsection{Descomposición de \texttt{periodo} en \texttt{anio} y \texttt{semestre}}

\begin{lstlisting}
df["periodo"] = df["periodo"].astype(str)
df["anio"]     = df["periodo"].str[:4].astype(int)
df["semestre"] = df["periodo"].str[4:]
\end{lstlisting}

La descomposición permite al modelo capturar:

\begin{itemize}
  \item \textbf{Tendencia temporal continua} a través de \texttt{anio} —
        efectos de política educativa, pandemia COVID-19 en 2020--2021.
  \item \textbf{Efecto de semestre} de forma independiente — el semestre 1 y
        el semestre 2 pueden tener perfiles de estudiantes distintos.
\end{itemize}

Mantener \texttt{periodo} como entero (20142, 20151\ldots) crearía una variable
con saltos irregulares que el modelo trataría incorrectamente como escala
numérica continua.

\subsection{Normalización de etiquetas categóricas}

\begin{lstlisting}
df["cole_area_ubicacion"] = df["cole_area_ubicacion"].str.strip().str.title()
df["cole_naturaleza"]     = df["cole_naturaleza"].str.strip().str.title()
df["estu_genero"]         = df["estu_genero"].str.strip().str.upper()
\end{lstlisting}

Los archivos crudos presentan inconsistencias de capitalización entre periodos
(\texttt{"URBANO"}, \texttt{"Urbano"}, \texttt{"urbano"}). La normalización es
necesaria antes de cualquier agrupación o filtro para evitar categorías duplicadas.

% ══════════════════════════════════════════════════════════════════════════════
\section{Segmentación por área: urbano / rural}
% ══════════════════════════════════════════════════════════════════════════════

\subsection{Criterio de segmentación}

Los subconjuntos \texttt{v2\_datos\_urbano.csv} y \texttt{v2\_datos\_rural.csv}
se generan filtrando por \texttt{cole\_area\_ubicacion} (área de ubicación del
\textbf{colegio}), no por la residencia del estudiante.

\textbf{Justificación:} el área del colegio define el contexto educativo que el
modelo intentará explicar: infraestructura disponible, cuerpo docente, jornada
escolar. La residencia del estudiante es relevante para acceso, pero muchos
estudiantes rurales asisten a colegios urbanos y viceversa. Usar el área del
colegio mantiene la coherencia entre el target (puntaje obtenido en ese colegio)
y el contexto institucional.

\subsection{Distribución resultante}

\vspace{4pt}
\begin{tabular}{lrr}
  \toprule
  \thead{Área} & \thead{Filas} & \thead{Porcentaje} \\
  \midrule
  \rowcolor{rowalt}
  Urbano         & 4\,185\,434 & 86.82\% \\
  Rural          &   635\,603  & 13.18\% \\
  \midrule
  \textbf{Total} & \textbf{4\,821\,037} & \textbf{100\%} \\
  \bottomrule
\end{tabular}

% ══════════════════════════════════════════════════════════════════════════════
\section{Impacto en memoria y relevancia para deep learning}
% ══════════════════════════════════════════════════════════════════════════════

\subsection{Reducción de peso en RAM}

\vspace{4pt}
\begin{tabular}{lcr}
  \toprule
  \thead{Dataset} & \thead{Columnas} & \thead{Peso en memoria} \\
  \midrule
  \rowcolor{rowalt}
  Original & 57  & 12\,706.9 MB ($\approx$12.4 GB) \\
  Limpio   & 28  & 7\,606.4 MB ($\approx$7.4 GB) \\
  \midrule
  \textbf{Reducción} & $\mathbf{-29}$ & \textbf{$-$5\,030.6 MB ($\approx$39.6\%)} \\
  \bottomrule
\end{tabular}

\subsection{Por qué importa en deep learning}

\begin{itemize}
  \item \textbf{Tiempo de entrenamiento.} Cada época procesa todos los registros.
        Cargar $\approx$5 GB menos por época reduce el tiempo de I/O y el
        movimiento de datos entre RAM y GPU/CPU de forma acumulativa a lo largo
        de cientos de épocas.

  \item \textbf{Presión sobre la VRAM.} Los batches se mueven a memoria de GPU.
        Un dataset más liviano permite usar \textit{batch sizes} más grandes, lo
        que estabiliza el gradiente y acelera la convergencia. Con 12 GB en RAM,
        mantener el dataset completo en VRAM de una GPU estándar (8--16 GB) era
        inviable; con 7.4 GB hay más margen.

  \item \textbf{Ruido en los gradientes.} Las columnas eliminadas (identificadores,
        flags, variables de baja varianza) no aportan señal pero sí ruido durante
        el \textit{backpropagation}. Cada columna irrelevante obliga a la red a
        aprender que su peso debe ser $\approx$0, desperdiciando capacidad del
        modelo y ralentizando la convergencia.

  \item \textbf{Embeddings categoriales.} Las redes neuronales representan
        variables categóricas como embeddings. Eliminar columnas de alta
        cardinalidad sin semántica útil (\texttt{cole\_nombre\_establecimiento}
        con miles de valores únicos) evita embeddings masivos que consumirían
        parámetros sin aportar poder predictivo.
\end{itemize}

% ══════════════════════════════════════════════════════════════════════════════
\section{Análisis exploratorio: decisiones metodológicas}
% ══════════════════════════════════════════════════════════════════════════════

\subsection{Filtro de edad al derivar \texttt{edad\_examen}}

La edad al momento del examen se calcula como:

\begin{lstlisting}
edad_tmp["fnac"]         = pd.to_datetime(
    edad_tmp["estu_fechanacimiento"], dayfirst=True, errors="coerce")
edad_tmp["edad_examen"]  = edad_tmp["anio"] - edad_tmp["fnac"].dt.year
edad_df = edad_tmp[edad_tmp["edad_examen"].between(14, 60)]
\end{lstlisting}

El campo \texttt{estu\_fechanacimiento} contiene errores de captura frecuentes
(años imposibles: 1900, 2090, fechas con formato inválido). El rango 14--60 cubre:

\begin{itemize}
  \item Mínimo realista para grado 11 (14 años).
  \item Máximo razonable para validantes de bachillerato adultos (grado 26).
\end{itemize}

\begin{nota}
  Este filtro es exclusivamente analítico (notebook exploratorio) y
  \textbf{no se aplica} al dataset de entrenamiento.
\end{nota}

\subsection{Codificación ordinal para correlación de Spearman}

Para calcular la correlación de Spearman entre variables categóricas y
\texttt{punt\_global}, se aplican mapeos ordinales manuales:

\vspace{4pt}
\begin{tabularx}{\textwidth}{>{\ttfamily}l X}
  \toprule
  \thead{Variable} & \thead{Escala} \\
  \midrule
  \rowcolor{rowalt}
  fami\_estratovivienda  & Sin Estrato=0, Estrato 1=1, \ldots, Estrato 6=6 \\
  cole\_area\_ubicacion  & Rural=0, Urbano=1 \\
  \rowcolor{rowalt}
  cole\_naturaleza       & Oficial=0, No Oficial=1 \\
  fami\_tieneX           & No=0, Si=1 \\
  \rowcolor{rowalt}
  fami\_educacionmadre/padre & Ninguno=0, \ldots, Postgrado=9 \\
  fami\_personashogar    & Una=1, Dos=2, \ldots, Diez O Más=10 \\
  \bottomrule
\end{tabularx}
\vspace{6pt}

Los mapeos respetan el orden natural de las categorías, haciendo la correlación
interpretable como ``a mayor nivel, mayor/menor puntaje''. Usar \textit{one-hot
encoding} para este análisis produciría correlaciones sin sentido direccional.

\begin{nota}
  Estos mapeos son solo para el análisis exploratorio. El modelo de deep
  learning usará embeddings o codificación específica para cada variable.
\end{nota}

% ══════════════════════════════════════════════════════════════════════════════
\end{document}
